% !TEX program = LuaLaTeX --shell-escape
\documentclass[10pt,a6paper]{article}
\pagestyle{empty}

% 1. MARGÓK BEÁLLÍTÁSA
\usepackage{geometry}
\geometry{
	top=1cm,
	bottom=1cm,
	left=1.5cm,
	right=1.5cm,
}

% Címsorok számozásának kikapcsolása
\setcounter{secnumdepth}{-1}

% betűtípus
\usepackage{fontspec}
\setmainfont{Linux Libertine O} % Linux Libertine betűtípus használata
\newfontfamily\litfont{liturgy.ttf}

% Bekezdések
\setlength{\parindent}{0pt}
\setlength{\parskip}{0.4em}

% Nyelv beállítása (latin)

\usepackage{polyglossia}
\setdefaultlanguage{latin} % Alapértelmezett nyelv latin
\setotherlanguage{hungarian} % Másodlagos nyelv magyar

% csomagok

\usepackage{microtype} %mikrotipografia
\usepackage{xcolor} %színek
\usepackage{lettrine} %inicialek

%==============================================================================
% Stilusok
%==============================================================================

% Focim
\newcommand{\titulus}[1]{%
  \par
  \vspace{1em}
  \begin{center}
    \Large\bfseries #1
  \end{center}
  \vspace{0.8em}
}

% Alcim
\newcommand{\pars}[1]{%
  \par
  \vspace{0.8em}
  \begin{center}
    \large\bfseries #1
  \end{center}
  \vspace{0.5em}
}

% Kisebb cim
\newcommand{\subpars}[1]{%
  \par
  \vspace{0.6em}
  {\bfseries #1\par}
  \vspace{0.3em}
}

% Rubrika
\newcommand{\rubrica}[1]{%
  \par
  {\color{red}\footnotesize #1\par}
  \smallskip
}

\newcommand{\V}{\textcolor{red}{\textbf{\litfont V.}}\ }
\newcommand{\R}{\textcolor{red}{\textbf{\litfont R.}}\ }

% Antifona
\newcommand{\antiphona}[1]{%
  \par
  \medskip
  \textbf{Ant.} #1\par
  \smallskip
}

% Oratio
\newcommand{\oratio}[1]{%
  \par
  \medskip
  {\bfseries Oratio}\par
  #1
  \medskip
}

% Zsoltar cime
\newcommand{\psalmus}[2]{%
  \par
  \vspace{0.8em}
  \begin{center}
    \bfseries Psalmus #1\\
    \itshape #2
  \end{center}
  \vspace{0.4em}
}

% Himnusz cime
\newcommand{\hymnus}[1]{%
  \par
  \vspace{0.8em}
  \begin{center}
    \bfseries Hymnus\\
    \itshape #1
  \end{center}
  \vspace{0.4em}
}

\newcommand{\forditas}[1]{%
  {\small\itshape (#1)}%
}

\newcommand{\initial}[2]{%
  \lettrine[
    lines=2,
    lhang=.1,
    findent=2pt,
    nindent=0pt
  ]{\color{red}\normalfont #1}{\normalfont #2}
}


\begin{document}

  \titulus{Ante Divinum Officum}
  
  \rubrica{
    flexis genibus
  }

  \initial{A}{peri} Dómine, os meum ad benedicéndum nomen sanctum tuum: 
  munda quoque cor meum ab ómnibus vanis, pervérsis et aliénis cogitatiónibus; 
  intelléctum illúmina, afféctum inflámma, ut digne, atténte ac devóte hoc Offícium recitáre váleam, 
  et exaudíri mérear ante conspéctum divínæ Maiestátis tuæ. Per Christum Dóminum nostrum. \R Amen
  \smallskip

  \forditas{
    Nyisd meg Uram, ajkamat neved dicséretére; tisztítsd meg szívemet minden hiábavaló 
    és vétkes gondolattól; világosítsd meg értelmemet, gyújtsd lángra szívemet, 
    hogy ezt a zsolozsmát méltón, figyelmesen, áhítattal végezzem, 
    és meghallgatásra méltó legyek isteni fölséged színe előtt. Krisztus a mi Urunk által.
  }
  \smallskip

  \initial{D}{ómine}, in unióne illíus divínæ intentiónis, qua ipse in terris laudes Deo persolvísti, 
  has tibi horas ({\color{red} \footnotesize vel} hanc tibi horam) persólvo.
  \smallskip

  \forditas{
    Uram, egyesülve azzal az isteni szándékkal, mellyel te magad teljesítetted 
    e földön élvén az Atyaisten dicséretét, akarom én is most ezt a zsolozsmát elvégezni.
  }

  \newpage
  
  \titulus{Post Divinum Officum}
  
  \rubrica{
    Dicitur autem flexis semper genibus in privata etiam recitatione, 
    præter quam ab iis, qui ob certam infirmitatem vel gravioris impedimenti causam nequeant genuflectere.
  }

  \initial{S}{acrosánctæ} et indivíduæ Trinitáti, crucifíxi Dómini nostri Iesu Christi humanitáti, 
  beatíssimæ et gloriosíssimæ sempérque Vírginis Maríæ fœcúndæ integritáti, 
  et ómnium Sanctórum universitáti sit sempitérna laus, honor, virtus et glória ab omni creatúra, 
  nobísque remíssio ómnium peccatórum, per infiníta sǽcula sæculórum. \R Amen
  \smallskip

  \forditas{
    A legszentebb és oszthatatlan Háromságnak, a mi megfeszített Urunk, Jézus Krisztus emberségének, 
    a Boldogságos és dicsőséges, mindenkor Szűz Mária termékeny érintetlenségének 
    és minden szentek egyességének legyen örök dicséret, tisztesség, erő és dicsőség minden 
    teremtménytől, nekünk pedig minden bűneink bocsánata örökkön örökké!
  }
  \smallskip

  \V Beáta víscera Maríæ Vírginis, quæ portavérunt ætérni Patris Fílium. \\
  \R Et beáta úbera, quæ lactavérunt Christum Dóminum. \\
  \smallskip

  \forditas{
    Boldog Szűz Mária méhe, mely az örök Atya Fiát hordozta. - 
    És boldogok az emlők, melyek Krisztust, az Urat szoptatták.
  }
  
  {\color{red} \footnotesize Et dicitur secreto} Pater noster. {\color{red} \footnotesize {et} Ave Maria.

\end{document}
